\subsection{星代数中的近似单位元}

定义 7. --- 设 A 是一个赋范代数。A 的一个近似单位元是 A 的单位球上的一个滤子基 F，使得对 A 中每个 x，A 上的滤子基 xF 与 Fx 收敛于 x，换言之：

\[
\lim _ {f, \mathfrak {F}} f x = \lim _ {f, \mathfrak {F}} x f = x.
\]

若 A 是一个星代数，则称近似单位元 F 是递增的，如果 F 是 \(A_{+}\) 上的一个滤子基。

设 A 是一个星代数。我们用 \(A_{+}^{\leqslant1}\)（相应地，\(A_{+}^{<1}\)）表示 A 中范数 \(\leqslant 1\)（相应地，范数 \(< 1\)）的正元素组成的集合；它们是 A 的谱包含在 \([0,1]\)（相应地，\([0,1)\)）中的 Hermite 元。

命题 17. --- 设 A 是一个星代数，\(\mathfrak{F}\) 是 \(\mathrm{A}_{+}^{\leqslant 1}\) 上的一个滤子基。为使 \(\mathfrak{F}\) 是 A 的一个递增近似单位元，必须且只需有

\[
\lim _ {f, \mathfrak {F}} f x = x
\]

对 A 的每个正元素 x 成立。

这个条件显然是必要的；我们来证明它是充分的。设 \(\widetilde{\mathrm{A}}\) 是由 A 添加单位元而得的星代数。设 \(x\) 是 A 的一个元素，\(f \in \mathrm{A}_{+}^{\leqslant 1}\)。我们有

\[
\begin{array}{c} \| f x - x \| ^ {2} = \| (f x - x) (f x - x) ^ {*} \| = \| (f - 1) x x ^ {*} (f - 1) \| \\ \leqslant \| (f - 1) x x ^ {*} \| \end{array}
\]

因为 \(\| f - 1\| \leqslant 1\)（I, p.~116 的引理 12, \(c\))）。于是我们有

\[
\limsup _ {f, \mathfrak {F}} \| f x - x \| ^ {2} \leqslant \limsup _ {f, \mathfrak {F}} \| f x x ^ {*} - x x ^ {*} \|.
\]

由于 \(xx^{*}\) 是正的 (I, p.~118, 定理 2)，假设蕴含

\[
\limsup _ {f, \mathfrak {F}} \| f x - x \| ^ {2} \leqslant \| x x ^ {*} - x x ^ {*} \| = 0,
\]

从而

\[
\lim _ {f, \mathfrak {F}} f x = x.
\]

由于 A 的对合是连续的且 \(\mathfrak{F}\) 是 \(A_{h}\) 上的一个滤子基，由此得

\[
\lim _ {f, \mathfrak {F}} x f = \lim _ {f, \mathfrak {F}} (f ^ {*} x ^ {*}) ^ {*} = \lim _ {f, \mathfrak {F}} (f x ^ {*}) ^ {*} = (x ^ {*}) ^ {*} = x.
\]

命题由此得证。

命题 18. --- 设 A 是一个星代数。有序集 \(A_{+}^{<1}\) 是向右定向的 (E, III, p.~12, 定义 7)，且它的截口滤子 (TG, I, p.~38, 例 2) 是 A 的一个递增近似单位元。

设 \(\widetilde{\mathsf{A}}\) 是由 A 添加记作 1 的单位元而得的星代数 (I, p.~106, 定义 4)。

设 \(g\) 是从 \([0,1)\) 到 \(\mathbf{R}_+\) 中的由 \(g(t) = t(1 - t)^{-1}\) 定义的函数。它是一个连续递增双射，其逆双射由 \(t \mapsto 1 - (1 + t)^{-1}\) 给出。

我们来证明有序集 \(A_{+}^{<1}\) 是向右定向的。设 x 与 y 是 \(A_{+}^{<1}\) 的元素。由于 \(g(0)=0\) 且 \(\mathrm{Sp}_{\mathrm{A}}^{\prime}(x)\) 与 \(\mathrm{Sp}_{\mathrm{A}}^{\prime}(y)\) 都包含在 \([0,1)\) 中，元素 \(g(x)\) 与 \(g(y)\) 有定义；它们是正的，于是 \(g(x)+g(y)\geqslant0\)。因此我们可以作 A 的元素 \(z=g^{-1}(g(x)+g(y))\)。我们有 \(\mathrm{Sp}_{\mathrm{A}}^{\prime}(z)\subset[0,1)\)，从而 \(z\in A_{+}^{<1}\)。

我们有 \(0 \leqslant g(x) \leqslant g(z)\)，于是 \(1 \leqslant 1 + g(x) \leqslant 1 + g(z)\)。I, p.~119 的引理 14 的断言 \(b\)) 蕴含 \(1 + g(x)\) 与 \(1 + g(z)\) 可逆，且 \((1 + g(z))^{-1} \leqslant (1 + g(x))^{-1}\)。从而 \(z = 1 - (1 + g(z))^{-1} \geqslant 1 - (1 + g(x))^{-1} = x\)。类似地，我们有 \(z \geqslant y\)。因此，\(z\) 优超 \(x\) 与 \(y\)（在 \(\mathrm{A}_{+}^{<1}\) 中）。于是有序集 \(\mathrm{A}_{+}^{<1}\) 是向右定向的。用 \(\mathfrak{F}\) 表示它的截口滤子

设 \(x\) 是 A 的一个正元素。对每个整数 \(n \geqslant 1\)，令 \(e_n = g^{-1}(nx)\)；我们有 \(e_n \in \mathrm{A}_+^{<1}\)。设 \(h_n\) 是 \(\mathbf{R}_+\) 上对一切 \(t \in \mathbf{R}_+\) 由 \(h_n(t) = t^2(1 - g^{-1}(nt)) = t^2 / (1 + nt)\) 定义的连续函数。我们有 \(|h_n(t)| \leqslant t / n\) 对一切 \(t \geqslant 0\) 成立，从而 \(\| x(1 - e_n)x \| = \| h_n(x) \| \leqslant \| x \| / n\)。特别地，\(x(1 - e_n)x\) 当 \(n\) 趋于无穷时趋于 0。

设 \(\varepsilon > 0\) 是一个实数。取整数 \(n\) 使 \(\|x(1 - e_n)x\| < \varepsilon\)。于是对每个 \(f \in \mathrm{A}_+^{<1}\)，只要 \(f \geqslant e_n\)，我们有

\[
\begin{array}{r l} & {\| x - f x \| ^ {2} = \| (1 - f) x \| ^ {2} = \| ((1 - f) x) ^ {*} (1 - f) x \| = \| x ^ {*} (1 - f) ^ {2} x \|} \\ & {\qquad \qquad \qquad \qquad \qquad \qquad \qquad \qquad \qquad \qquad = \| x (1 - f) ^ {2} x \|.} \end{array}
\]

此外，由于 \(0 \leqslant f \leqslant 1\)，由此得 \((1-f)-(1-f)^{2} = (1-f)f \geqslant 0\)（I, p.~116 的引理 12, d)），从而 \((1-f)^{2} \leqslant 1-f\)。由于 \(1-f \leqslant 1-e_{n}\)，由此得 \(0 \leqslant (1-f)^{2} \leqslant 1-e_{n}\)。依 I, p.~119 的引理 14, a) 与 c)，由此得

\[
\left\| x - f x \right\| ^ {2} = \left\| x (1 - f) ^ {2} x \right\| \leqslant \left\| x (1 - e _ {n}) x \right\| <   \varepsilon .
\]

于是 \(\lim_{f,\mathfrak{F}}fx = x\) 对一切 \(x \in A_{+}\) 成立。因此依命题 17，滤子 \(\mathfrak{F}\) 是 A 的一个近似单位元。

